% ============================================================
% REWRITE DRAFT — Chapter 2, sections 4 (Results) and 5 (Conclusion).
% Economics-paper register. Facts about the measures only; no method.
% NOTE the espoused-enacted gap figures are RECOMPUTED on the rebuilt
% espoused scores (+1.84 warmth, +0.88 strictness); the old +0.51 / +0.07
% predate the rebuild and are retired.
% ============================================================

\section{Results}
\label{sec:p1_results}

This section reports what the measures show. Three results stand out: the
enacted and espoused measures diverge, and they diverge more for warmth than
for strictness; every public source of text tracks the espoused warmth measure
and none tracks the enacted one; and behaviour policy text is, if anything,
inversely related to observed practice.

\textbf{Distributions.} \Cref{tab:score_dist} reports the measures by sample
and \cref{fig:score_distributions} their distributions. Enacted warmth across
the \NTierOne{} visited schools has mean 6.6 and standard deviation 1.0 on the
0--10 scale; enacted strictness has mean 7.0 and standard deviation 0.9.
Espoused scores across the interviewed schools are wider on warmth (SD 1.5)
and comparable on strictness (SD 0.9). The two enacted dimensions are
correlated at $r = 0.59$: warm schools tend also to be orderly ones, which is
consistent with the authoritative-school framework and inconsistent with a
single warm-to-strict axis. Between-school variation is genuinely narrow in
the English system, and score compression partly reflects that: about
two-thirds of secondary schools are graded Good and a further sixth
Outstanding, and common inspection, curriculum and training pressures
compress the distribution of any culture measure.

\textbf{Enacted and espoused culture diverge.} On the \NTierOne{} schools
with both, headteachers rate their schools warmer than observers record by
$1.8$ points on the 0--10 scale and stricter by $0.9$ (paired
$p < 0.001$ for both). The two sources are correlated at $\GoldWarmthSplit$
for warmth and $\GoldStrictnessSplit$ for strictness. The gap is therefore
larger, and the agreement lower, for warmth: what a headteacher says about
relationships is a poorer guide to what a visitor finds than what they say
about behaviour. Within the interview, the items on which heads report
reliably are those about the consistency of practice (sanctions applied
consistently, $r = 0.39$ with observed strictness; rewards applied
consistently, $0.29$), not those about how people feel (pupils feel safe,
$0.16$; teachers care, $0.11$, neither significant). Headteachers report
accurately on whether their systems are applied and inaccurately on
relational climate.

\textbf{Public text tracks talk, not practice.} \Cref{tab:source_criterion}
sets each source's warmth and strictness score against the enacted and the
espoused criterion. The pattern is uniform. For strictness, the
inspection-report score is correlated with the observed measure at
$\CorrOfstedStrictness$ and with the espoused measure at
$\CorrOfstedStrictEspoused$; behaviour policies and websites carry no
strictness signal against either. For warmth, no source is correlated with
the observed measure, and the inspection report and the parent survey are
correlated with the espoused one ($0.29$ and $0.31$). The parent survey is
the fourth independent respondent group and shows the same crossover: parents'
reports of their children being happy and safe track what heads say about
warmth ($0.31$) more than what observers record ($0.16$), and track observed
strictness ($0.28$) better than observed warmth. \Cref{fig:crossover_warmth}
draws the warmth half of the table. What is written and reported about a
school's warmth, by inspectors, by parents and by the school itself, tracks
the school's account of itself and not the classroom.

\begin{table}[htbp]
\centering
\small
\caption{Each text source's warmth and strictness score correlated with the
enacted (observed) and espoused (headteacher-reported) gold-standard measure.
Enacted correlations use the visited schools; espoused correlations use the
interviewed schools. The parent survey has no strictness item; its warmth
composite is also shown against enacted strictness.}
\label{tab:source_criterion}
\begin{tabular}{lcccc}
\toprule
 & \multicolumn{2}{c}{\textit{Warmth score vs}} & \multicolumn{2}{c}{\textit{Strictness score vs}} \\
\cmidrule(lr){2-3}\cmidrule(lr){4-5}
\textit{Source} & enacted & espoused & enacted & espoused \\
 & $(n \approx 101)$ & $(n \approx 299)$ & $(n \approx 101)$ & $(n \approx 299)$ \\
\midrule
Ofsted inspection report & $-0.07$ & $+0.29$ & $+0.40$ & $+0.38$ \\
Behaviour policy         & $-0.08$ & $+0.00$ & $+0.01$ & $+0.00$ \\
School website           & $-0.03$ & $-0.07$ & $+0.09$ & $+0.07$ \\
Parent survey            & $+0.16$ & $+0.31$ & \multicolumn{2}{c}{($+0.28$ vs enacted strictness)} \\
\bottomrule
\end{tabular}
\begin{minipage}{0.86\linewidth}
\vspace{4pt}\scriptsize
\textit{Notes}: Pearson correlations. Text scores are the instruments of
record in \cref{tab:instruments_adopted}. Enacted and espoused measures are
defined in Section~3. Correlations of $|r| \geq 0.20$ are significant at the
five per cent level at these sample sizes.
\end{minipage}
\end{table}


\textbf{Behaviour policies describe aspiration.} The behaviour-policy warmth
score is uncorrelated with observed warmth ($-0.08$) and negatively
correlated with the parent measure ($-0.04$, $p = 0.03$); the strictness score
is uncorrelated with observed strictness ($0.01$). One school in ten publishes
a policy authored by its trust rather than by the school. Written behaviour
policies are governance documents that state intent, and schools that state
the warmest intent are not, on these measures, the warmest schools. The same
result appears in the following chapter as a small negative association
between behaviour-policy warmth and pupil progress.

\textbf{What the text does and does not contain.} Two further results
qualify the inspection-report scores. The warmth score tracks the inspection
grade at $-0.52$ from report text stripped of every stated grade, where human
raters of the same text track it at $-0.25$: inspectors write about
relationships in a register that moves with their overall verdict, and the
model reads that register. The score is therefore a characterisation of the
report rather than a clean measure of the school. The strictness score also
tracks the verdict ($-0.51$) but, unlike warmth, it tracks observed practice
as well, and in the following chapter its association with pupil progress
survives conditioning on the grade. Finally, the one website construct with an
external criterion, the prominence of religious faith, is recovered from
website text with an area under the ROC curve of 0.960: the method reads what
is there, and warmth is not.

\begin{figure}[htbp]
  \centering
  \includegraphics[width=0.85\textwidth]{fig_crossover_warmth.pdf}
  \caption{Each source's warmth score correlated with the espoused criterion
  (filled) and the enacted criterion (open). Every source tracks what
  headteachers report better than what observers record.}
  \label{fig:crossover_warmth}
\end{figure}

% ============================================================
\section{Conclusion}
\label{sec:p1_conclusion}

This paper constructs measures of warmth and strictness for English secondary
schools from direct observation, headteacher interview and public text, and
asks how far each dimension can be recovered from data that exist for every
school. Three findings summarise it.

Strictness is legible in public text and warmth is not. An inspection
report's strictness score is correlated with observed strictness at
$\CorrOfstedStrictness$, and a text-only model predicts it out of sample at
$\RidgeBStrictness$; no text source is correlated with observed warmth, and
the text-only warmth model has no out-of-sample skill. The asymmetry is a
fact about the sources, not the instruments: inspectors record conduct and
enforcement in detail and describe relationships in a uniform positive
register, schools' own documents describe intent, and websites are not written
about the thing observers see.

Espoused and enacted culture diverge, and more for warmth. Headteachers
over-state warmth relative to observation by nearly two points and agree with
observers only weakly on it, while their reports of behaviour systems being
applied consistently track observed strictness. Every public source, including
a survey of parents, tracks the espoused warmth measure and not the enacted
one. The gap is dimension-specific and appears in four independent kinds of
evidence, which suggests it is a feature of how schools present themselves
rather than an artefact of any one instrument.

The instruments carry a measurement result of their own. Across roughly
85 versions of eleven language-model instruments, agreement with reference
labels was unrelated to correlation with the observed measure, and the
best-agreeing version was usually the worse measure; instruments that ask for
facts outperformed those that ask for judgements; scoring related constructs
in separate calls produced correlations between them that the observations
do not show; and an unchanged rubric moved between models lost most of its
agreement. These results apply to any study that scores text with a language
model against human labels.

The measures have limits that the following chapter inherits. The visited
sample is 103 schools, weighted toward London and toward stronger schools; the
enacted scores come from one visit each; the inspection-report scores are read
from prose that carries the inspection verdict; and the reference labels for
the text instruments were produced by a language model rather than by human
raters, with the observed measures as the final check. Within those limits,
the enacted scores are a reliable, two-dimensional measure of school culture
for the visited schools, and the inspection-report strictness score extends
one of the two dimensions to the national population.
